\section{Introducción a la Modulación Digital}

Hasta ahora estudiamos sistemas con:
\begin{itemize}
    \item Mensaje analógico $x(t)$;
    \item $x_c(t)$ que depende en forma continua del mensaje.
\end{itemize}

Pasamos ahora a sistemas que comunican una sucesión de símbolos
$\{a_k\}_{k\in\mathbb{Z}}, a_k\in \mathcal{A}$, alfabeto finito.
Caso típico $\mathcal{A}=\{0,1\}$, comunicación binaria. También
se usan sistemas M-arios, con $M$ símbolos.

\paragraph{Motivación}
\begin{enumerate}
    \item Símbolos codifican muchas cosas (texto, imágenes, o señales analógicas muestreadas y cuantificadas).
    \item En un canal con distorsión y/o ruido, recuperar una señal analógica implica deshacer la distorsión. En el caso digital, solo hace falta una detección de qué símbolo se transmitió.
\end{enumerate}


\begin{ejemplo}$\phantom{1}$

    \FloatBarrier

    \begin{figure}[htb]
        \centering\begin{minipage}{\textwidth}
        \bpgf{Imagenes/modDig/modDig-1}
        \begin{tikzpicture}
            \draw (0,0) -- (1,0) -- (1,1) -- (2,1) node [anchor=south west] {$x(t)$} -- (2,0) -- (3,0);
            \begin{scope}[shift={(5,0)}]
                \draw (0,0) -- (5,0);
                \draw (0,1) -- (5,1);
                \path (2.5,0.5) node {CANAL};
            \end{scope}
            \begin{scope}[shift={(12,0)}]
                \draw (-0.5,0) -- (0,0);
                \draw plot [domain=0:1,samples=200] (\x,{sin(pi*\x r)+0.1*rand});
                \draw (1,0) -- (1.5,0) node [anchor=south] {$y(t)$};
            \end{scope}
        \end{tikzpicture}
        \endpgfgraphicnamed\end{minipage}
    \end{figure}
    \FloatBarrier

    $y(t)$ llega muy distorsionada, pero a partir de ella puedo decidir si el pulso es hacia arriba o hacia abajo, en forma relativamente sencilla (en comparación con lo que sería recomponer una forma de onda analógica). En particular, esto permite regenerar el pulso.

\end{ejemplo}


    No es que la comunicación digital esté libre de errores; si
    pretendemos transmisión libre de errores, hay límites
    intrínsecos de información que puede pasar por el canal.
    La ventaja es de \emph{implementación}, trabajar
    con símbolos discretos abre más posibilidades. La digitalización
    llegó a todo (telecomunicaciones y electrónica) en forma
    paralela.
%
%     al igual que lo que ocurre al
%    comparar electrónica digital vs. analógica. La a todo en
%    forma paralela.

%\end{remarkSN}

Aún si transmitimos símbolos, la comunicación digital interactúa
con aspectos analógicos a dos niveles:


\begin{itemize}
    \item los símbolos pueden venir de una conversión A/D de un mensaje analógico.
    \item para enviar los símbolos por un canal analógico se deben representar
    por formas de onda. Un caso típico: tren de pulsos $\displaystyle \sum_k a_k p(t-kT_s)$ ``PAM'': \emph{pulse amplitude modulation}.
\end{itemize}

\begin{figure}[htb]
    \centering\begin{minipage}{\textwidth}
    \bpgf{Imagenes/modDig/modDig-2}
    \begin{blockdiagram}
        \matrix[row sep=5mm, column sep=10mm]{
        & \node[etiqueta] (inicio) {$x(t)$}; &\\
        \node[etiqueta] (otro) {\begin{minipage}{3cm}
                        otras fuentes de símbolos (texto, imágenes, ...)
                    \end{minipage}}; & \node[block] (conv) {A/D}; &\\
        & \node[block] (mod) {MOD. DIGITAL}; &\\
        & & \node [block] (canal) {CANAL};\\
        };
        \draw [->] (inicio) -- (conv);
        \draw [->] (conv) -- (mod);
        \draw [->] (otro) |- (mod);
        \draw [->] (mod) |- node [anchor=north] {$\sum a_k p(t-kT_s)$} (canal);
    \end{blockdiagram}
    \endpgfgraphicnamed\end{minipage}
\end{figure}

La idea de PCM (pulse-code-modulation, enviar mensajes analógicos
codificados en trenes de pulsos) data de 1937.  Hoy en día, es más
natural pensar en la comunicación de símbolos como la operación
básica, independiente del origen de los símbolos.

Nos concentraremos primero entonces en la capa inferior del
diagrama: cómo enviar símbolos discretos por un canal.

\subsection*{Algunos parámetros}
\begin{itemize}
    \item Velocidad de señalización, $f_s$. Se mide en baudios $=
    \rfrac{\text{símbolos}}{seg}$.
    \item Velocidad de transmisión de información, $vti$. Se mide en $bps$.

    Si tengo $M$ símbolos en el alfabeto, cada uno tiene $\log_2(M)$ bits de información.
    \[\Rightarrow vti= f_s \log_2(M).\]
    \item Ancho de banda, $B$ del canal en $Hz$.
    \item Eficiencia espectral en $\rfrac{bps}{Hz}$:
    $\frac{vti}{B}$.
\end{itemize}

O sea, $vti$ es cómo se ve la velocidad en la capa superior
(ej.: ADSL $512\,Kbps$). Esa velocidad, como veremos, dependerá
del ancho de banda físico, $B$.

La relación entre ellos depende de cómo se elijan las formas de onda que representan los símbolos. Dos temas:
\begin{itemize}
    \item la forma de los pulsos.
    \item como se mapean los símbolos en pulsos: ``codificación de línea''.
\end{itemize}
